\documentclass[11pt]{article}

\usepackage[a4paper,margin=1in]{geometry}
\usepackage{amsmath,amssymb,amsthm,mathtools}
\usepackage{graphicx}
\usepackage{float}
\usepackage{subcaption}
\usepackage{enumitem}
\usepackage{microtype}
\usepackage{xcolor}
\usepackage{hyperref}
\usepackage{fancyhdr}
\usepackage{lastpage}

\hypersetup{
  colorlinks=true,
  linkcolor=blue,
  urlcolor=blue,
  citecolor=blue
}

\setlist[enumerate]{itemsep=0.25em,topsep=0.25em}
\setlength{\parindent}{0pt}
\setlength{\parskip}{0.5em}

\theoremstyle{definition}
\newtheorem{exercise}{Exercise}[section]
\newenvironment{solution}{\begin{proof}[Solution]}{\end{proof}}

\DeclareMathOperator{\Area}{Area}
\DeclareMathOperator{\Vol}{Vol}

\newcommand{\HKUSTCrest}{\detokenize{../Week 2/hkust-crest.png}}

\pagestyle{fancy}
\fancyhf{}
\setlength{\headheight}{18.1pt}
\renewcommand{\headrulewidth}{0.4pt}
\fancyhead[L]{\IfFileExists{\HKUSTCrest}{\includegraphics[height=14pt]{\HKUSTCrest}}{}}
\fancyhead[C]{MATH1014 Calculus II Tutorial}
\fancyhead[R]{Week 5}
\fancyfoot[C]{\thepage\ /\ \pageref{LastPage}}

\title{\vspace{-1.0em}MATH1014 Calculus II Tutorial\\Week 5}
\author{Wenkai Fan}
\date{}

\begin{document}
\maketitle
\thispagestyle{fancy}

\section{Review}

\subsection{Integration by Parts}
\begin{enumerate}
  \item Formula:
  \[
  \int u\,dv = uv-\int v\,du.
  \]
  \item A practical rule of thumb for choosing $u$: pick something that becomes simpler after differentiating (e.g.\ $\ln x$, $\arcsin x$, powers of $\ln x$).
  \item If integration by parts produces the original integral again, treat it as an unknown and solve a linear equation for it (this happens often for $\sin(\ln x)$ and $\cos(\ln x)$ type integrals).
\end{enumerate}

\subsection{Trigonometric Integrals and Substitutions}
\begin{enumerate}
  \item Use the derivative cycle $\frac{d}{dx}(\sin x)=\cos x$ and $\frac{d}{dx}(\cos x)=-\sin x$. This lets you build a numerator as the derivative of a convenient denominator, e.g.\ $(\sin x+\cos x)'=\cos x-\sin x$.
  \item Useful identities:
  \[
  \cos(\alpha-\beta)=\cos\alpha\cos\beta+\sin\alpha\sin\beta,
  \qquad
  \sin^2 x+\cos^2 x=1.
  \]
  \item Tangent half-angle substitution $u=\tan\frac{x}{2}$:
  \[
  \sin x=\frac{2u}{1+u^2},\quad
  \cos x=\frac{1-u^2}{1+u^2},\quad
  dx=\frac{2}{1+u^2}\,du.
  \]
\end{enumerate}

\section{Exercises}

\subsection{Integration by Parts}

\begin{exercise}[Q1]
Evaluate $\displaystyle \int \frac{\arcsin x}{\sqrt{1-x}}\,dx$.
\end{exercise}

\begin{solution}
\textbf{Key idea.} Write the integral in a form suitable for integration by parts using $d(\sqrt{1-x})$.
Since $d(\sqrt{1-x})=-\frac{1}{2\sqrt{1-x}}\,dx$, we have
\[
\int \frac{\arcsin x}{\sqrt{1-x}}\,dx=-2\int \arcsin x\,d(\sqrt{1-x}).
\]
Integration by parts gives
\begin{align*}
-2\int \arcsin x\,d(\sqrt{1-x})
&=-2\sqrt{1-x}\,\arcsin x+2\int \sqrt{1-x}\cdot \frac{d}{dx}(\arcsin x)\,dx \\
&=-2\sqrt{1-x}\,\arcsin x+2\int \frac{\sqrt{1-x}}{\sqrt{1-x^2}}\,dx.
\end{align*}
Since $\sqrt{1-x^2}=\sqrt{(1-x)(1+x)}$, the remaining integrand simplifies to $\frac{1}{\sqrt{1+x}}$:
\[
2\int \frac{\sqrt{1-x}}{\sqrt{(1-x)(1+x)}}\,dx
=2\int \frac{1}{\sqrt{1+x}}\,dx
=4\sqrt{1+x}+C.
\]
Therefore,
\[
\int \frac{\arcsin x}{\sqrt{1-x}}\,dx=-2\sqrt{1-x}\,\arcsin x+4\sqrt{1+x}+C.
\]
\end{solution}

\begin{exercise}[Q2]
Evaluate $\displaystyle \int x^2\ln x\,dx$.
\end{exercise}

\begin{solution}
Use integration by parts with $u=\ln x$ and $dv=x^2\,dx$. Then $du=\frac{1}{x}dx$ and $v=\frac{x^3}{3}$.
Hence
\[
\int x^2\ln x\,dx=\frac{x^3}{3}\ln x-\int \frac{x^3}{3}\cdot \frac{1}{x}\,dx
=\frac{x^3}{3}\ln x-\frac{1}{3}\int x^2\,dx
=\frac{x^3}{3}\ln x-\frac{x^3}{9}+C.
\]
\end{solution}

\begin{exercise}[Q3]
Evaluate $\displaystyle \int \ln^2 x\,dx$.
\end{exercise}

\begin{solution}
Use integration by parts with $u=\ln^2 x$ and $dv=dx$. Then $du=\frac{2\ln x}{x}dx$ and $v=x$.
So
\[
\int \ln^2 x\,dx=x\ln^2 x-2\int \ln x\,dx.
\]
Apply integration by parts again for $\int \ln x\,dx$ with $u=\ln x$, $dv=dx$:
\[
\int \ln x\,dx=x\ln x-\int 1\,dx=x\ln x-x+C.
\]
Substitute back:
\[
\int \ln^2 x\,dx=x\ln^2 x-2(x\ln x-x)+C
=x\ln^2 x-2x\ln x+2x+C.
\]
\end{solution}

\begin{exercise}[Q4]
Evaluate $\displaystyle \int \cos(\ln x)\,dx$.
\end{exercise}

\begin{solution}
Let
\[
I=\int \cos(\ln x)\,dx,\qquad J=\int \sin(\ln x)\,dx.
\]
Integration by parts for $I$ with $u=\cos(\ln x)$, $dv=dx$ gives $du=-\frac{\sin(\ln x)}{x}dx$ and $v=x$, so
\[
I=x\cos(\ln x)+\int \sin(\ln x)\,dx=x\cos(\ln x)+J.
\]
Similarly for $J$ with $u=\sin(\ln x)$, $dv=dx$ gives $du=\frac{\cos(\ln x)}{x}dx$ and $v=x$, hence
\[
J=x\sin(\ln x)-\int \cos(\ln x)\,dx=x\sin(\ln x)-I.
\]
Substitute this into $I=x\cos(\ln x)+J$:
\[
I=x\cos(\ln x)+x\sin(\ln x)-I
\quad\Rightarrow\quad
2I=x\bigl(\cos(\ln x)+\sin(\ln x)\bigr).
\]
Therefore
\[
\int \cos(\ln x)\,dx=\frac{x}{2}\bigl(\cos(\ln x)+\sin(\ln x)\bigr)+C.
\]
\end{solution}

\subsection{Trigonometric Integrals and Substitutions}

\begin{exercise}[Q5]
Evaluate $\displaystyle \int \frac{\cos x}{\sin x+\cos x}\,dx$.
\end{exercise}

\begin{solution}
Rewrite the numerator:
\[
\cos x=\frac{1}{2}\bigl(\sin x+\cos x\bigr)+\frac{1}{2}\bigl(\cos x-\sin x\bigr).
\]
Hence
\begin{align*}
\int \frac{\cos x}{\sin x+\cos x}\,dx
&=\frac{1}{2}\int 1\,dx+\frac{1}{2}\int \frac{\cos x-\sin x}{\sin x+\cos x}\,dx \\
&=\frac{x}{2}+\frac{1}{2}\ln|\sin x+\cos x|+C.
\end{align*}
\end{solution}

\begin{exercise}[Q6]
Evaluate $\displaystyle \int \frac{1}{\sin(x+a)\cos(x+b)}\,dx$.
\end{exercise}

\begin{solution}
Use $\cos(a-b)=\cos(x+a)\cos(x+b)+\sin(x+a)\sin(x+b)$:
\begin{align*}
\int \frac{1}{\sin(x+a)\cos(x+b)}\,dx
&=\frac{1}{\cos(a-b)}\int \frac{\cos(a-b)}{\sin(x+a)\cos(x+b)}\,dx \\
&=\frac{1}{\cos(a-b)}\int\left(\frac{\cos(x+a)}{\sin(x+a)}+\frac{\sin(x+b)}{\cos(x+b)}\right)dx \\
&=\frac{1}{\cos(a-b)}\left(\ln|\sin(x+a)|-\ln|\cos(x+b)|\right)+C \\
&=\frac{1}{\cos(a-b)}\ln\left|\frac{\sin(x+a)}{\cos(x+b)}\right|+C.
\end{align*}
\end{solution}

\begin{exercise}[Q7]
Evaluate $\displaystyle \int \frac{1}{(2+\cos x)\sin x}\,dx$.
\end{exercise}

\begin{solution}
Multiply by $\dfrac{\sin x}{\sin x}$ to create $d(\cos x)$:
\begin{align*}
\int \frac{1}{(2+\cos x)\sin x}\,dx
&=\int \frac{\sin x\,dx}{(2+\cos x)\sin^2 x}
=-\int \frac{d(\cos x)}{(2+\cos x)(1-\cos^2 x)}.
\end{align*}
Let $t=\cos x$. Then $dt=-\sin x\,dx$ and
\[
\int \frac{1}{(2+\cos x)\sin x}\,dx
=-\int \frac{1}{(2+t)(1-t^2)}\,dt.
\]
Use the identity $3=(2+t)+(1-t)$ to split the integrand:
\begin{align*}
-\int \frac{1}{(2+t)(1-t^2)}\,dt
&=-\frac{1}{3}\int \frac{(2+t)+(1-t)}{(2+t)(1-t^2)}\,dt \\
&=-\frac{1}{3}\int \frac{1}{1-t^2}\,dt-\frac{1}{3}\int \frac{1}{(2+t)(1+t)}\,dt.
\end{align*}
Now decompose each part:
\[
\frac{1}{1-t^2}=\frac{1}{2}\left(\frac{1}{1-t}+\frac{1}{1+t}\right),
\qquad
\frac{1}{(2+t)(1+t)}=-\frac{1}{2+t}+\frac{1}{1+t}.
\]
Therefore,
\begin{align*}
\int \frac{1}{(2+\cos x)\sin x}\,dx
&=-\frac{1}{6}\int\left(\frac{1}{1-t}+\frac{1}{1+t}\right)dt
-\frac{1}{3}\int\left(-\frac{1}{2+t}+\frac{1}{1+t}\right)dt \\
&=-\frac{1}{6}\bigl(-\ln|1-t|+\ln|1+t|\bigr)
-\frac{1}{3}\bigl(-\ln|2+t|+\ln|1+t|\bigr)+C \\
&=\frac{1}{6}\ln|1-t|+\frac{1}{3}\ln|2+t|-\frac{1}{2}\ln|1+t|+C.
\end{align*}
Substitute $t=\cos x$ and combine logs:
\[
\int \frac{1}{(2+\cos x)\sin x}\,dx
=\frac{1}{6}\ln\left(\frac{(1-\cos x)(2+\cos x)^2}{(1+\cos x)^3}\right)+C.
\]
\end{solution}

\begin{exercise}[Q8]
Evaluate $\displaystyle \int \frac{1}{2\sin x-\cos x+5}\,dx$.
\end{exercise}

\begin{solution}
Use the tangent half-angle substitution $u=\tan\frac{x}{2}$, so
\[
\sin x=\frac{2u}{1+u^2},\qquad \cos x=\frac{1-u^2}{1+u^2},\qquad dx=\frac{2}{1+u^2}\,du.
\]
Then
\begin{align*}
\int \frac{1}{2\sin x-\cos x+5}\,dx
&=\int \frac{\frac{2}{1+u^2}}{\frac{4u}{1+u^2}-\frac{1-u^2}{1+u^2}+5}\,du
=\int \frac{1}{3u^2+2u+2}\,du \\
&=\frac{1}{3}\int \frac{1}{\left(u+\frac{1}{3}\right)^2+\frac{5}{9}}\,du
=\frac{1}{\sqrt{5}}\arctan\left(\frac{3u+1}{\sqrt{5}}\right)+C \\
&=\frac{1}{\sqrt{5}}\arctan\left(\frac{3\tan\frac{x}{2}+1}{\sqrt{5}}\right)+C.
\end{align*}
\end{solution}

\end{document}
